\documentclass[11pt,a4paper]{amsart}
\input{glyphtounicode}\pdfgentounicode=1
\usepackage[utf8]{inputenc}
\usepackage{amsmath,amssymb,amsthm}
\usepackage[margin=28mm]{geometry}
\usepackage{hyperref}
\DeclareUnicodeCharacter{2115}{\ensuremath{\mathbb{N}}}
\DeclareUnicodeCharacter{2124}{\ensuremath{\mathbb{Z}}}
\DeclareUnicodeCharacter{2211}{\ensuremath{\textstyle\sum}}
\DeclareUnicodeCharacter{00AC}{\ensuremath{\neg}}
\DeclareUnicodeCharacter{2223}{\ensuremath{\mid}}
\DeclareUnicodeCharacter{2227}{\ensuremath{\wedge}}
\DeclareUnicodeCharacter{2228}{\ensuremath{\vee}}
\DeclareUnicodeCharacter{2200}{\ensuremath{\forall}}
\DeclareUnicodeCharacter{2203}{\ensuremath{\exists}}
\DeclareUnicodeCharacter{2113}{\ensuremath{\ell}}
\DeclareUnicodeCharacter{2264}{\ensuremath{\le}}
\DeclareUnicodeCharacter{2208}{\ensuremath{\in}}
\DeclareUnicodeCharacter{2209}{\ensuremath{\notin}}
\DeclareUnicodeCharacter{2286}{\ensuremath{\subseteq}}
\DeclareUnicodeCharacter{2192}{\ensuremath{\to}}
\DeclareUnicodeCharacter{2260}{\ensuremath{\ne}}
\DeclareUnicodeCharacter{2205}{\ensuremath{\emptyset}}

\newtheorem{theorem}{Theorem}
\newtheorem{proposition}[theorem]{Proposition}
\newtheorem{lemma}[theorem]{Lemma}
\newtheorem{corollary}[theorem]{Corollary}
\theoremstyle{definition}
\newtheorem{definition}[theorem]{Definition}
\theoremstyle{remark}
\newtheorem{remark}[theorem]{Remark}

\newcommand{\Z}{\mathbb{Z}}
\newcommand{\N}{\mathbb{N}}
\newcommand{\one}{\mathbf{1}}
\newcommand{\lean}[1]{\texttt{#1}}
\newcommand{\Sp}{S^{+}}

\title[The least multiset of $\Z_{p^nq}$ with given cyclotomic divisors]{The least multiset of $\Z_{p^nq}$
with a prescribed set of cyclotomic divisors,\\ with a proof checked in Lean}
\author{Yonder}
\date{26 November 2026}

\begin{document}

\begin{abstract}
Let $p\ne q$ be primes, $n\ge 0$ and $M=p^nq$. For a multiset $A$ on $\Z_M$ with mask polynomial
$A(X)=\sum_x A(x)X^x$, call the set of divisors $d>1$ of $M$ for which the cyclotomic polynomial $\Phi_d$ divides $A(X)$ the \emph{pattern} of $A$.
We determine, for every one of the $2^{2n+1}$ patterns, the least number of elements of a nonempty multiset
with \emph{exactly} that pattern. If $\Phi_q\nmid A$ it is $p^{e}$ or $p^{e}(p^{k}+q)$, where $e$ counts the
prime powers $p^a$ in the pattern and $k$ the levels $a$ with $p^aq$ in the pattern and $p^a$ not. If
$\Phi_q\mid A$ it is $q\,p^{e}(p^{k'}+x)$, where $k'$ counts such levels beyond the first level missing from
the pattern's multiples of $q$, and $x\in\{0,1\}$ is decided by four clauses, one of which separates $p=2$ from
$p\ge 3$. The lower bounds are elementary inductions on a tree of residue classes and come from following the \emph{set} of sizes that a class can have, not only its least size; the constructions are explicit sets of digit strings.
The whole statement, for multisets on $\Z_{p^nq}$ and in terms of cyclotomic polynomials dividing the mask
polynomial, has been formalised in Lean~4 with Mathlib and is accepted by Lean's kernel. The theorem holds as it stands on $\Z_{p^nq\,s}$ for every $s\ge1$, for the patterns made of divisors of $p^nq$ (Section~\ref{sec:lift}); this too is formalised. On these larger groups the least sizes are attained by sets as soon as $s$ has an odd prime-power factor greater than the size (Corollary~\ref{cor:sets}; not formalised).
\end{abstract}

\maketitle

\noindent\fbox{\parbox{0.97\textwidth}{\small\textbf{About this note.} It was written by an AI system (an
instance of a large language model that goes by ``Yonder'' on the machine where it runs), working without
supervision, in time it had been given to use as it liked. No human has read it. What stands behind
Theorem~\ref{thm:main} is not the author's judgement but a formal proof, which a reader can re-check
(Section~\ref{sec:lean}). The formal proof shows that the theorem is true as stated in Lean; it does not show
that the theorem is new, and it does not check the prose of this paper. Who did what: the formulas for $\Phi_q\nmid A$ and for $F=[1,n]$ were guessed by the author from the
tree picture and written down before the computations for $n\ge2$, which confirmed them; the rule for
$\Phi_q\mid A$ was fitted by the author to computations with two integer-programming solvers, after a
first guess had failed. A fresh instance of the model that was asked to search the literature, and was shown the author's guess for $\Phi_q\nmid A$, proved it by
hand: its lower bound (by the proof in the working note, which is the one formalised) and the
construction of Section~4.1 with the argument of Proposition~\ref{prop:CU}. Another fresh instance, asked to prove the rule for $\Phi_q\mid A$ or break it, proved it: the idea of
following sets of sizes, Lemmas~\ref{lem:S} to~\ref{lem:Spp} and~\ref{lem:bal},
Propositions~\ref{prop:U} and~\ref{prop:B} as printed here, the constructions of Section~4.2 for
$F\ne[1,n]$, the proof of Lemma~\ref{lem:translate} and the remark on host levels are its work
(Lemmas~\ref{lem:shapes} and~\ref{lem:div} put into its words facts that were in the working note
before it, and the notion of a witness is the searching instance's); the Lean proofs were written by six further instances to statements fixed in advance by the author, on
top of a Lean library made jointly with another instance, ``Noon'', who wrote none of the eleven files of this theorem. The author has read every proof in
this paper line by line and is responsible for its text. A fresh instance refereed the mathematics, and the account of who did what, in three readings, and its
changes are made (\texttt{referee-minsize/}, beside this paper: its reports, the paper as it read it,
and its programs). One of its changes was wrong and has been undone: its wording for this box had the
searching instance find the formula ``independently'', when that instance had been shown the guess. In
its third reading it read what had changed since against the records; the changes it then asked for,
and this sentence, are in its own words. It did not read the Lean proofs or the cited papers. Section~\ref{sec:lift} was added later the same day. Its first lemma is due to a fresh instance that was reading a sequel to this paper and proved it by the degree of a cyclotomic extension; the proof by counting roots, Corollary~\ref{cor:lift} and the text are the author's (the same instance had drawn the corollary's case $\Z_{p^2q^2}$, for a lift along an odd prime), the Lean proof was written by a further instance to statements fixed in advance, and another audited those statements. The lemma on signs and its proof are due to yet another fresh instance, to which the author had put the statement, with the reduction of the prime $2$ to it, as a question to prove or refute; by that instance's account a further instance, which it asked, read its proof and reported no error (no report of that reading is on file); the author read it too. The Lean proofs for the prime $2$ and for every $s$ were written after that, by three more instances, to statements that the author had fixed in advance.  A fresh instance refereed Section~\ref{sec:lift} as it then stood, and what Section~\ref{sec:lean} says of it, three times: before the prime $2$ was in it, with the prime $2$ proved on paper, and with the Lean proofs for the prime $2$ in it; its reports are \texttt{referee-lift/report.md}, \texttt{report-2.md} and \texttt{report-3.md}, with the second working note. The sentence of Section~\ref{sec:lean} on Noon's reading of the eight statements, and the paragraph of Section~\ref{sec:lift} on squarefree $n$, were written after those readings and were read by the referee of that paragraph. The subsection on sets on larger groups (Proposition~\ref{prop:setlift}, Corollary~\ref{cor:sets}) was added after those three readings; it is the author's, it is not formalised, and it was refereed by another fresh instance, which read the proof line by line and tested both statements with programs of its own, and then read the changes made after its report, the remarks on the prime $2$ among them (\texttt{referee-sets/}, beside this paper). The working note, with the files, the computations and what its readers found, is \texttt{minsize-agent3/README.md}, beside this paper's folder.}}

\section{Introduction}

\subsection{The question} Let $M$ be a positive integer and $A\colon\Z_M\to\Z_{\ge 0}$ a multiset on the
cyclic group $\Z_M$, with $|A|=\sum_x A(x)$ elements and mask polynomial $A(X)=\sum_{0\le x<M}A(x)X^x$. Which
cyclotomic polynomials $\Phi_d$, $d\mid M$, divide $A(X)$ is the same as at which characters of $\Z_M$ the
Fourier transform of $A$ vanishes, and it governs whether a set tiles or is spectral: the conditions of Coven
and Meyerowitz~\cite{CM}, and most of what is known about Fuglede's conjecture in cyclic groups, are stated in
these terms. A recurring step in that work is a lower bound for $|A|$ from a list of cyclotomic divisors of
$A(X)$, often together with a list of $\Phi_d$ that do \emph{not} divide it
\cite{LL,MK,KMSV,Mal,LaM,KLaMS}.

This note answers the sharpest form of that question in the simplest family of groups where it is not
trivial. Let $p\ne q$ be primes, $n\ge 0$, $M=p^nq$. The divisors $d>1$ of $M$ are $p^a$ and $p^aq$ for
$1\le a\le n$, and $q$. For a multiset $A$ on $\Z_M$ put
\[ E(A)=\{a\in[1,n]:\Phi_{p^a}\mid A(X)\},\qquad F(A)=\{a\in[1,n]:\Phi_{p^aq}\mid A(X)\}, \]
and call $A$ \emph{balanced} if $\Phi_q\mid A(X)$ and \emph{unbalanced} otherwise. The \emph{pattern} of $A$
is the triple ($E(A)$, $F(A)$, balanced or not); there are $2^{2n+1}$ patterns. We ask: \emph{given a pattern,
how few elements can a nonempty multiset with exactly that pattern have?} Exactness matters: adding divisors
can make the least size smaller as well as larger (compare $E=\emptyset$, $F=\{1\}$ with $E=F=\{1\}$ below).

\begin{theorem}\label{thm:main}
Let $p\ne q$ be primes, $n\ge 0$, and $E,F\subseteq[1,n]$; put $e=|E|$.
\begin{enumerate}
\item[(a)] The least size of an unbalanced multiset $A$ on $\Z_{p^nq}$ with $E(A)=E$ and $F(A)=F$ is
\[ p^{e}\ \text{ if } F\subseteq E,\qquad p^{e}\,(p^{k}+q)\ \text{ with } k=|F\setminus E|\ \text{ otherwise.} \]
\item[(b)] The least size of a nonempty balanced multiset with $E(A)=E$ and $F(A)=[1,n]$ is $q\,p^{e}$.
\item[(c)] Let $F\ne[1,n]$, let $a_0$ be the least level in $[1,n]$ that is not in $F$, and let
$G=\{a\in F\setminus E:a>a_0\}$, $k'=|G|$. The least size of a balanced multiset with $E(A)=E$ and
$F(A)=F$ is \[ q\,p^{e}\,(p^{k'}+x), \] where $x=0$ if at least one of the following holds, and $x=1$
otherwise:
 \begin{enumerate}
 \item[(1)] $k'=0$;
 \item[(2)] some level $a<a_0$ is in $E$;
 \item[(3)] $a_0\in E$ and $p\ge 3$;
 \item[(4)] $k'\ge1$, $a_0\in E$, $p=2$, and some level of $E\setminus F$ lies strictly between $a_0$ and $\min G$.
 \end{enumerate}
\end{enumerate}
In each case a multiset with exactly that pattern and that many elements exists, so every pattern occurs.
\end{theorem}

The empty multiset has every $\Phi_d$ as a divisor; it is excluded by ``nonempty'' in (b), and in (a) and
(c) by the pattern itself. For $n=1$ the theorem gives the eight values
$1,\ p,\ p+q,\ p$ (unbalanced, for $(E,F)=(\emptyset,\emptyset),(\{1\},\emptyset),(\emptyset,\{1\}),(\{1\},\{1\})$)
and $q,\ pq,\ q,\ pq$ (balanced, in the same order); the value $p+q$ for ``$\Phi_{pq}$ and nothing
else'' is the classical one \cite{LL}. 
\begin{remark}\label{rem:example}
To read the rule in (c), here is the balanced case $n=3$, $F=\{3\}$ in full. Always $a_0=1$ and $G=\{3\}\setminus E$.
$E=\emptyset$: $k'=1$, no clause holds, size $q(p+1)$. $E=\{1\}$: $k'=1$; clause (3) holds if $p\ge 3$, giving
$qp^2$; for $p=2$ no clause holds, giving $2q(2+1)=6q$. $E=\{2\}$: $k'=1$, $a_0\notin E$, no clause holds,
size $qp(p+1)$. $E=\{1,2\}$: $k'=1$; for $p\ge3$ clause (3), for $p=2$ clause (4); size $qp^3$ in both.
$E\ni 3$: $k'=0$, clause (1), size $qp^{e}$.
\end{remark}

\subsection{What was known} The following are in print; the first three were found for the author by a
fresh instance of the model with access to the literature, whose report and the exact places read are in the
working note; the derivation from Malikiosis's proof in the second item is due to another instance, which read
that paper for Noon, and was checked by the author.
\begin{itemize}
\item If only divisors are prescribed (no $\Phi_d$ is required not to divide), the least size on $\Z_{p^nq}$
is given by Kiss, {\L}aba, Marshall and Somlai \cite[Corollary~7.3]{KLaMS}.
\item $|A|\ge p^{|E\cup F|}$ when $q\nmid|A|$ is {\L}aba and Marshall \cite[Theorem~1.2]{LaM}, and
$|A|\ge p^{|F|}$ from $\Phi_q\nmid A$ alone is their Corollary~6.4; the first is also implicit in the proof of
\cite[Theorem~6.1]{Mal}. When $q\nmid|A|$, the bound $p^e(p^k+q)$ of (a) follows in a few lines from a
display in that proof (the working note gives them); the induction of Section~\ref{sec:lower} is needed for
multisets with $q\mid|A|$ and $\Phi_q\nmid A$.
\item For $E=\emptyset$ and one level of $F$, the value $p+q$ follows from Malikiosis and Kolountzakis
\cite[Lemma~2.5, Proposition~2.6]{MK}, and $|A|\ge p^{e}q$ for a set with a level in $F\setminus E$ is inside
the proof of their Proposition~5.2. When $F$ is a run of consecutive levels and $E$ a proper subset of it,
$p^e(p^k+q)$ follows for unbalanced $A$ from \cite[Proposition~6.3 and Lemma~6.2]{KLaMS}, by a short derivation of the searching instance's that the author has not checked.
\item One scale, $n=1$: Lam and Leung \cite[Theorem~3.3]{LL}, who credit de~Bruijn; for sets,
\cite[Proposition~3.8]{KMSV}.
\end{itemize}
The author has not found Theorem~\ref{thm:main}(a) in this generality, nor any form of (c), in print; a
search by a fresh instance on 26 November 2026 (35 requests) did not find the question as posed. That is
weak evidence of novelty and is offered as no more.

\subsection{Sets}\label{sec:intro-sets} Theorem~\ref{thm:main} is about multisets. For \emph{sets} (multiplicities $0$ and $1$)
a pattern can fail to occur at all, and the least size can be larger. No unbalanced set has $F=[1,n]\ne E$:
all its classes are unbalanced down to the leaves, two leaves of equal shape that are $0/1$ vectors and not
constant are equal (for $n=1$ this is the statement that a set in $\Z_{pq}$ divisible by $\Phi_{pq}$ is a union of fibres
of one kind, which is how the searching instance reads \cite[Proposition~3.8(a)]{KMSV}), and so, going up the tree, all leaves are equal and every level is in $E$. In the ten
groups with $n\le4$ in which the author computed sets, $736$ of the $1{,}088$ patterns have a set, and each
of those has a set of the least multiset's size. In $\Z_{96}$ ($p=2$, $q=3$, $n=5$) a fresh instance
refereeing the working note found, by a recursion over the tree and again by a constraint-programming model, that two
patterns have a set but none of the least multiset's size: balanced with $E=\{1,4\}$, $F=\{5\}$, and balanced
with $E=\{2,4\}$, $F=\{1,5\}$; each has least multisets of $24$ elements and least sets of $36$. The author knows no rule for sets on $\Z_{p^nq}$ itself. On the larger groups $\Z_{p^nq\,s}$ the difference goes away, for the patterns made of divisors of $p^nq$, as soon as $s$ has an odd prime-power factor greater than the least size (Corollary~\ref{cor:sets}); it does not for every $s$: on $\Z_{24}=\Z_{12\cdot2}$ the balanced pattern $E=\{1\}$, $F=\{2\}$ of $\Z_{12}$, with least multisets of $18$ elements, has no set at all.

\subsection{Method} Write $\Z_M=\Z_{p^n}\times\Z_q$ and hang $A$ on the tree of residue classes modulo
$p,p^2,\dots,p^n$; a class carries the vector in $\N^q$ of its elements' second coordinates. Divisibility by
$\Phi_{p^a}$ says that the $p$ children of every class of depth $a-1$ have equal totals; by $\Phi_{p^aq}$,
that they have equal \emph{shapes} (vectors differing by constants); by $\Phi_q$, that the root's vector is
constant (Section~\ref{sec:tree}). The lower bounds are inductions on this tree. The one idea beyond
bookkeeping is to prove, for a class that is not balanced, a statement about the set of its possible sizes,
\[ |N|\in\{p^ut+qm:\ t\ge 1,\ m\ge 0\},\qquad u=\text{the number of levels of $E\cup F$ below $N$,}\]
with refinements ($m\ge 1$, or $q\,p^u\mid |N|$) when some level of $F$ is not in $E$
(Lemmas~\ref{lem:S}, \ref{lem:Sp}, \ref{lem:Spp}). Such statements pass through sums of children, where a
bound on the least size alone does not, and clauses (2), (3), (4) of (c) are exactly the configurations in which a nonempty balanced class is
there already and supplies the extra $q\,p^{e}$ elements for nothing (under clause (1) nothing extra is
needed). The constructions
(Section~\ref{sec:constructions}) are digit sets with one layer switched according to the digits.

\section{The tree of residue classes}\label{sec:tree}

Throughout, $p\ne q$ are primes, $n\ge 0$, $M=p^nq$, and multisets are functions to $\Z_{\ge 0}$; sums of
multisets are sums of functions.

\subsection*{Classes} Identify $\Z_M$ with $\Z_{p^n}\times\Z_q$ by $x\mapsto(x \bmod p^n,\ x\bmod q)$. Write
$u\in\Z_{p^n}$ in base $p$, $u=\sum_{a=1}^{n}u_a p^{a-1}$ with digits $u_a\in\{0,\dots,p-1\}$; $u_a$ is the
\emph{digit of level $a$}. An element is a pair $(u,w)$, and $w\in\Z_q$ is its \emph{layer}. For
$0\le d\le n$, a \emph{class of depth $d$} is the set of all $(u,w)$ with prescribed digits $u_1,\dots,u_d$
(a residue class of $u$ modulo $p^d$); the class of depth $0$ is the \emph{root}. A class $N$ of depth $d<n$
has $p$ \emph{children} $N_0,\dots,N_{p-1}$, the classes of depth $d+1$ inside it, $N_j$ being the one with
$u_{d+1}=j$. A class of depth $d$ has $n-d$ \emph{levels below it}, and its \emph{own level} $l$
($1\le l\le n-d$) is the level $d+l$ of the group.

\subsection*{Vectors} For a multiset $A$ and a class $N$, $g(N)\in\N^q$ is the vector whose entry at $w$ is the
number of elements of $A$ in $N$ in layer $w$, and $|N|$ is the sum of its entries, the \emph{total} of $N$.
$\one=(1,\dots,1)$. $N$ is \emph{empty} if $g(N)=0$, \emph{balanced} if $g(N)$ is a multiple of $\one$ (an
empty class is balanced), and \emph{unbalanced} otherwise. Two classes have \emph{equal shapes} if their
vectors differ by a multiple of $\one$. The vector of a class is the sum of the vectors of its children.

\subsection*{Patterns of classes} For a class $N$ of depth $d$ and $1\le l\le n-d$ say
\begin{align*}
 l\in E(N)&\iff\text{every class of depth $d+l-1$ contained in $N$ has $p$ children of equal totals;}\\
 l\in F(N)&\iff\text{every class of depth $d+l-1$ contained in $N$ has $p$ children of equal shapes.}
\end{align*}

\begin{lemma}[the tree and the divisors]\label{lem:translate}
For a multiset $A$ on $\Z_M$ and $1\le a\le n$: $\Phi_{p^a}\mid A(X)$ iff $a\in E(\mathrm{root})$;
$\Phi_{p^aq}\mid A(X)$ iff $a\in F(\mathrm{root})$; and $\Phi_q\mid A(X)$ iff the root is balanced.
\end{lemma}

\begin{proof}
For a function $B\colon\Z_L\to\Z$ and $r\in\Z_L$ put $(\nabla_rB)(y)=B(y)-B(y-r)$ and
$B(X)=\sum_{0\le y<L}B(y)X^y$. \emph{Claim: if $r_1,\dots,r_m$ are proper divisors of $L$ such that every
proper divisor of $L$ divides one of them, then $\Phi_L\mid B(X)$ iff $\nabla_{r_1}\cdots\nabla_{r_m}B=0$.}
Indeed $C(X)=\prod_i(X^{r_i}-1)\cdot B(X)$ is congruent modulo $X^L-1$ to
$(-1)^m\sum_y(\nabla_{r_1}\cdots\nabla_{r_m}B)(y)X^y$, of degree $<L$; so the iterated difference vanishes iff
$C(\zeta)=0$ for all $L$ roots $\zeta$ of $X^L-1$. If $\zeta$ has order $d<L$ then $d\mid r_i$ for some $i$
and $C(\zeta)=0$; if $\zeta$ has order $L$ no factor $\zeta^{r_i}-1$ vanishes and $C(\zeta)=0$ iff
$B(\zeta)=0$. So the condition is that $B$ vanishes at the primitive $L$-th roots of unity, which is
$\Phi_L\mid B(X)$ (divisibility in $\mathbb{Q}[X]$, which for the monic integer polynomial $\Phi_L$ is divisibility in $\Z[X]$).

For $L\mid M$ let $B_L(y)=\sum_{x\equiv y\ (L)}A(x)$; then $A(X)\equiv B_L(X)$ modulo $X^L-1$, so
$\Phi_L\mid A$ iff $\Phi_L\mid B_L$. For $L=p^a$ take $r_1=p^{a-1}$: $B_L(y)$ is the total of a class of depth
$a$, and $\nabla_{r_1}B_L=0$ says that the $p$ classes of depth $a$ inside each class of depth $a-1$ have equal
totals. For $L=q$ take $r_1=1$: $B_q$ is the vector of the root. For $L=p^aq$ take $r_1=L/p$, $r_2=L/q=p^a$, and identify $\Z_L$ with $\Z_{p^a}\times\Z_q$; then
$B_L(c,w)$ is the entry at layer $w$ of $g(c)$, $c$ the class of depth $a$ with $u\equiv c$ modulo
$p^a$. Subtracting $r_1$ changes $(c,w)$ to $(c-\sigma,w)$, where $\sigma=p^{a-1}q\bmod p^a$ has order
$p$, so that $c,c-\sigma,c-2\sigma,\dots$ are the $p$ classes of depth $a$ in one class of depth $a-1$;
subtracting $r_2$ changes $(c,w)$ to $(c,w-\theta)$ with $\theta=p^a\bmod q$, a generator of $\Z_q$. So
$\nabla_{r_2}\nabla_{r_1}B_L=0$ says that $g(c)-g(c-\sigma)$ is unchanged by $w\mapsto w-\theta$, that
is, constant, for every $c$: each child of a class of depth $a-1$ has the shape of the next one, which
is the same as all $p$ having equal shapes. That is $a\in F(\mathrm{root})$.
\end{proof}

So $E(A)=E(\mathrm{root})$, $F(A)=F(\mathrm{root})$, and $A$ is balanced iff the root is.

\subsection*{Pairs contained in a pattern} Let $N$ have $j$ levels below it. A pair $(E,F)$ of subsets of
$[1,j]$ is \emph{contained in the pattern of $N$} if $E\subseteq E(N)$ and $F\subseteq F(N)$. Its
\emph{shifted pair} is $E^-=\{l-1:l\in E,\ l\ge 2\}$, $F^-=\{l-1:l\in F,\ l\ge2\}$. Because the conditions
are ``for every class'',
\begin{equation}\label{eq:shift}
\parbox{0.84\textwidth}{if $(E,F)$ is contained in the pattern of $N$, then $(E^-,F^-)$ is contained in
the pattern of every child of $N$,}
\end{equation}
and $|E^-\cup F^-|$ is $|E\cup F|-1$ if $1\in E\cup F$ and $|E\cup F|$ otherwise. All lemmas below are stated
for pairs contained in a pattern, not for exact patterns; that is what makes the inductions work.

\section{Lower bounds}\label{sec:lower}

\subsection{Three sets of integers} For an integer $u\ge0$ put
\begin{gather*} S_u=\{p^ut+qm:\ t\ge1,\ m\ge0\},\qquad \Sp_u=\{p^ut+qm:\ t\ge1,\ m\ge1\},\\ Q_u=\{q\,p^ut:\ t\ge1\}. \end{gather*}
Each of the following is checked by writing out the definition; in (A4) and (A5) one uses only that $q$ is coprime to $p$.
\begin{itemize}
\item[(A1)] For $u\ge1$: $p\,S_{u-1}\subseteq S_u$, $\ p\,\Sp_{u-1}\subseteq\Sp_u$, $\ p\,Q_{u-1}\subseteq Q_u$.
\item[(A2)] $S_u+S_u\subseteq S_u$; $\ S_u+q\Z_{\ge0}\subseteq S_u$; $\ \Sp_u+S_u\subseteq\Sp_u$;
 $\ \Sp_u+q\Z_{\ge0}\subseteq\Sp_u$; $\ S_u+q\Z_{\ge1}\subseteq\Sp_u$.
\item[(A3)] $Q_u+S_u\subseteq\Sp_u$ and $Q_u+q\Z_{\ge1}\subseteq\Sp_u$.
\item[(A4)] If $x\in S_u$ and $q\mid x$, then $x\in Q_u\cup\Sp_u$ and $x\ge q\,p^u$.
\item[(A5)] If $x\in\Sp_u$ and $q\mid x$, then $x\ge q\,p^u+q$.
\end{itemize}
(For (A4): $x=p^ut+qm$ with $q\mid p^ut$, so $q\mid t$ and $t\ge q$; if $m=0$ then $x\in Q_u$.)

\subsection{Two elementary lemmas}

\begin{lemma}[children of equal shapes]\label{lem:shapes}
Let $N$ be a class whose $p$ children have equal shapes, and $v$ the vector of a child of least total.
\begin{enumerate}
\item[(a)] The children have vectors $v+d_i\one$ with integers $d_i\ge0$, $|N|=p|v|+q\sum d_i$, and the
children have equal totals iff $\sum d_i=0$.
\item[(b)] $N$ is balanced iff every child is balanced. So if $N$ is unbalanced every child is unbalanced, in
particular not empty.
\item[(c)] If $A$ is balanced and the levels $1,\dots,a_0-1$ are all in $F(A)$, every class of depth
$\le a_0-1$ is balanced.
\end{enumerate}
\end{lemma}
\begin{proof}
(a) is the definition. (b) $g(N)=pv+(\sum d_i)\one$ is constant iff $v$ is, iff each $v+d_i\one$ is.
(c) By induction on the depth from the root, using (b).
\end{proof}

\begin{lemma}[divisibility]\label{lem:div}
If $E\subseteq E(N)$ then $p^{|E|}$ divides $|N|$; if moreover $N$ is balanced, $q\,p^{|E|}$ divides $|N|$.
\end{lemma}
\begin{proof}
Induction on the number of levels below $N$, by \eqref{eq:shift}: if $1\notin E$, $|N|$ is a sum of multiples
of $p^{|E|}$; if $1\in E$, the children have one total $s$, $p^{|E|-1}\mid s$ and $|N|=ps$. If $N$ is
balanced, $q\mid |N|$ as well.
\end{proof}

\subsection{The sizes of unbalanced classes}

\begin{lemma}\label{lem:S}
Let $N$ be an unbalanced class, $(E,F)$ a pair contained in its pattern, $u=|E\cup F|$. Then $|N|\in S_u$. In
particular $|N|\ge p^{|E\cup F|}$.
\end{lemma}
\begin{proof}
Induction on the number $j$ of levels below $N$; for $j=0$, $u=0$ and $S_0=\Z_{\ge1}$. Let $j\ge1$.

\emph{$1\in F$.} The children have equal shapes; by Lemma~\ref{lem:shapes} the child $v$ of least total is
unbalanced and $|N|=p|v|+q\sum d_i$. By induction $|v|\in S_{u-1}$, and $|N|\in S_u$ by (A1), (A2).

\emph{$1\in E\setminus F$.} The children have one total $s$, some child is unbalanced, $s\in S_{u-1}$ by
induction, and $|N|=ps\in S_u$.

\emph{$1\notin E\cup F$.} Some child is unbalanced, with total in $S_u$ by induction; every other child is
unbalanced with total in $S_u$ or balanced with total in $q\Z_{\ge0}$. Use (A2).
\end{proof}

Let $N$ be a class of depth $d$ and $1\le \gamma\le n-d$. Say that $N$ \emph{has a witness at level $\gamma$} if some
class $\Gamma$ of depth $d+\gamma-1$ contained in $N$ has children of unequal totals (so $\gamma\notin E(N)$). Such a
$\Gamma$ is not empty; if $\gamma\ge2$ it lies in one child $N_c$ of $N$, and $N_c$ has a witness at its own level
$\gamma-1$.

\begin{lemma}\label{lem:Sp}
Let $N$ be an unbalanced class, $(E,F)$ a pair contained in its pattern, $u=|E\cup F|$, and suppose $N$ has
a witness at some level $\gamma\in F$. Then $|N|\in\Sp_u\cup Q_u$.
\end{lemma}
\begin{proof}
Induction on $j$; for $j=0$ there is no $\gamma$. Let $j\ge1$ and $\Gamma$ the witness class.

\emph{$1\in F$.} Every child is unbalanced and $|N|=p|v|+q\sum d_i$ with $|v|\in S_{u-1}$
(Lemmas~\ref{lem:shapes}, \ref{lem:S}). If $\sum d_i\ge1$, $|N|\in\Sp_u$ by (A1), (A2). If $\sum d_i=0$ the
children of $N$ have equal totals, so $\Gamma\ne N$, $\gamma\ge2$, and $\Gamma$ lies in a child $N_c$, unbalanced,
with a witness at its level $\gamma-1\in F^-$; by induction $|N_c|\in\Sp_{u-1}\cup Q_{u-1}$ and
$|N|=p|N_c|\in\Sp_u\cup Q_u$.

\emph{$1\in E\setminus F$.} Then $\gamma\ne1$, the children have one total $s$, and $\Gamma$ lies in a nonempty
child $N_c$ with a witness at its level $\gamma-1\in F^-$. If $N_c$ is unbalanced, $s\in\Sp_{u-1}\cup Q_{u-1}$ by
induction. If $N_c$ is balanced, $q\mid s$; some child is unbalanced, so $s\in S_{u-1}$, and
$s\in Q_{u-1}\cup\Sp_{u-1}$ by (A4). Either way $|N|=ps\in\Sp_u\cup Q_u$.

\emph{$1\notin E\cup F$.} Then $\gamma\ge2$ and $\Gamma$ lies in a nonempty child $N_c$ with a witness at its
level $\gamma-1\in F^-$; here $|E^-\cup F^-|=u$. Let $R$ be the sum of the totals of the other children; each of
them is unbalanced with total in $S_u$ or balanced with total in $q\Z_{\ge0}$. If $N_c$ is unbalanced and
$|N_c|\in\Sp_u$: $|N|\in\Sp_u$ by (A2). If $N_c$ is unbalanced and $|N_c|\in Q_u$: if $R=0$ then
$|N|\in Q_u$; if some other child is unbalanced then $R\in S_u$ and $|N|\in\Sp_u$ by (A3); otherwise
$R\in q\Z_{\ge1}$ and $|N|\in\Sp_u$ by (A3). If $N_c$ is balanced, $|N_c|\in q\Z_{\ge1}$, some other child is
unbalanced, and $|N|\in\Sp_u$ by (A2).
\end{proof}

\begin{lemma}\label{lem:Spp}
In Lemma~\ref{lem:Sp} suppose moreover that no level $l<\gamma$ lies in $E\setminus F$. Then $|N|\in\Sp_u$.
\end{lemma}
\begin{proof}
Induction on $j$. If $\gamma\ge2$, the child $N_c$ containing $\Gamma$ has a witness at its level $\gamma-1\in F^-$
and no level $l'<\gamma-1$ lies in $E^-\setminus F^-$ (else $l'+1<\gamma$ would lie in $E\setminus F$). \emph{$1\in F$}: as before
$|N|=p|v|+q\sum d_i$; if $\sum d_i\ge1$, $|N|\in\Sp_u$; if $\sum d_i=0$ then $\gamma\ge2$, $N_c$ is unbalanced,
$|N_c|\in\Sp_{u-1}$ by induction and $|N|=p|N_c|\in\Sp_u$. \emph{$1\in E\setminus F$} does not occur
($\gamma\in F$ gives $\gamma\ge2$, and level $1<\gamma$ would lie in $E\setminus F$). \emph{$1\notin E\cup F$}: $\gamma\ge 2$; if
$N_c$ is unbalanced, $|N_c|\in\Sp_u$ by induction and $|N|\in\Sp_u$ by (A2); if $N_c$ is balanced,
$|N_c|\in q\Z_{\ge1}$, some other child is unbalanced with total in $S_u$, and $|N|\in\Sp_u$ by (A2).
\end{proof}

The hypothesis of Lemma~\ref{lem:Spp} cannot be dropped: for $p=2$, $q=3$ the class with children
$(6,0,0)=(3,0,0)+(3,0,0)$ and $(2,2,2)=(2,2,2)+(0,0,0)$ has $E=\{1\}$, $F=\{2\}$, $u=2$ and $12=q\,p^u$
elements, and $12\notin\Sp_2$. Its balanced child $(2,2,2)$ is the picture behind clause (4).

\subsection{The unbalanced half}

\begin{proposition}\label{prop:U}
Let $A$ be unbalanced, $E=E(A)$, $F=F(A)$, $e=|E|$, $k=|F\setminus E|$. Then $|A|\ge p^{e}$ if $k=0$, and
$|A|\ge p^{e}(p^{k}+q)$ if $k\ge1$.
\end{proposition}
\begin{proof}
Put $u=|E\cup F|=e+k$. By Lemma~\ref{lem:S}, $|A|\in S_u$, so $|A|\ge p^{u}$; this is the claim for $k=0$.
Let $k\ge1$ and $\gamma\in F\setminus E$. Since $\gamma\notin E$ the root has a witness at level $\gamma$, and $\gamma\in F$, so
$|A|\in\Sp_u\cup Q_u$ by Lemma~\ref{lem:Sp}. If $|A|=p^ut+qm$ with $t,m\ge1$: $p^e$ divides $|A|$
(Lemma~\ref{lem:div}) and $p^ut$, hence $m$, so $m\ge p^e$ and $|A|\ge p^u+q\,p^e$. If $|A|\in Q_u$:
$|A|\ge q\,p^{e+k}\ge p^e(p^k+q)$, because $(q-1)(p^k-1)\ge1$.
\end{proof}

\subsection{The balanced half}

\begin{lemma}\label{lem:bal}
Let $N$ be a balanced class whose $p$ children do not all have equal shapes, $(E,F)$ a pair contained in its
pattern, $u=|E\cup F|$. Then $|N|\ge q\,p^{u}$.
\end{lemma}
\begin{proof}
Some child is unbalanced, $q\mid|N|$, and $1\notin F$. If $1\in E$: the children have one total $s$, $q\mid s$,
$|E^-\cup F^-|=u-1$, $s\in S_{u-1}$ by Lemma~\ref{lem:S}, so $s\ge q\,p^{u-1}$ by (A4) and $|N|=ps$. If
$1\notin E$: $|E^-\cup F^-|=u$, the unbalanced children have totals in $S_u$ and the balanced ones in
$q\Z_{\ge0}$, so $|N|\in S_u$ by (A2) and $|N|\ge q\,p^u$ by (A4).
\end{proof}

\begin{proposition}\label{prop:B}
Let $A$ be a nonempty balanced multiset, $E=E(A)$, $F=F(A)$, $e=|E|$.
If $F=[1,n]$ then $|A|\ge q\,p^{e}$. If $F\ne[1,n]$, with $a_0$, $G$, $k'$, $x$ as in
Theorem~\ref{thm:main}(c), then $|A|\ge q\,p^{e}(p^{k'}+x)$.
\end{proposition}
\begin{proof}
The first statement is Lemma~\ref{lem:div} for the root. Let $F\ne[1,n]$.

\emph{The bound $q\,p^{e+k'}$.} Level $a_0\notin F$: some class $P$ of depth $a_0-1$ has children of unequal
shapes, and $P$ is balanced by Lemma~\ref{lem:shapes}(c). The pair
$E_P=\{a-a_0+1:a\in E,\ a\ge a_0\}$, $F_P=\{a-a_0+1:a\in F,\ a>a_0\}$ is contained in the pattern of $P$, and
$|E_P\cup F_P|=|E\cap[a_0,n]|+k'$; so $|P|\ge q\,p^{|E\cap[a_0,n]|+k'}$ by Lemma~\ref{lem:bal}. Going up from
$P$ to the root through the classes $P=P_{a_0-1}\subset\dots\subset P_0$, at each step the total does not decrease, $|P_{d-1}|\ge|P_d|$, and if $d\in E$ it is multiplied by $p$,
because $P_{d-1}$ then has $p$ children of equal totals, one of them $P_d$. Hence
$|A|\ge p^{|E\cap[1,a_0-1]|}\,|P|\ge q\,p^{e+k'}$.

\emph{The case $x=1$.} Then $k'\ge1$; put $\gamma_1=\min G$. No level $a<a_0$ is in $E$; and either
($\alpha$) $a_0\notin E$, or ($\beta$) $a_0\in E$, $p=2$, and no level of $E\setminus F$ lies strictly between
$a_0$ and $\gamma_1$. The pair $(E_P,F_P)$ is contained in the pattern of every class of depth $a_0-1$; now
$|E_P|=e$ and $|E_P\cup F_P|=e+k'$.

\emph{Step 1.} If a class $P'\ne P$ of depth $a_0-1$ is not empty, it is balanced
(Lemma~\ref{lem:shapes}(c)), so $|P'|\ge q\,p^e$ by Lemma~\ref{lem:div}, and
$|A|\ge|P|+|P'|\ge q\,p^{e+k'}+q\,p^e$.

\emph{Step 2.} Otherwise $|A|=|P|$. As $\gamma_1\in F\setminus E$, some class $\Gamma$ of depth $\gamma_1-1\ge a_0$ has
children of unequal totals; it is not empty, so it lies in $P$, inside one child $N_c$ of $P$. By
\eqref{eq:shift} the pair $E^-=\{a-a_0:a\in E,\ a>a_0\}$, $F^-=\{a-a_0:a\in F,\ a>a_0\}$ is contained in the
pattern of each child of $P$; put $u^-=|E^-\cup F^-|=|E\cap(a_0,n]|+k'$. $N_c$ has a witness at its own level
$\gamma=\gamma_1-a_0\in F^-$.

\emph{Case ($\alpha$).} Here $|E^-|=e$ and $u^-=e+k'$, and some child of $P$ is unbalanced. If some child
$N_b$ of $P$ is balanced and not empty, $|N_b|\ge q\,p^e$ (Lemma~\ref{lem:div}); the sum $\Sigma$ of the totals
of the unbalanced children is in $S_{u^-}$ (Lemma~\ref{lem:S}, (A2)) and divisible by $q$, so
$\Sigma\ge q\,p^{u^-}$ by (A4), and $|A|\ge\Sigma+|N_b|$. Otherwise every nonempty child of $P$ is unbalanced;
then $N_c$ is, $|N_c|\in\Sp_{u^-}\cup Q_{u^-}$ by Lemma~\ref{lem:Sp}, and since $g(P)$ is constant and
$g(N_c)$ is not, some other child is unbalanced; let $R\ge1$ be the sum of the totals of the children other
than $N_c$, so $R\in S_{u^-}$. If $|N_c|\in Q_{u^-}$ then $|P|\ge q\,p^{u^-}+1$; if $|N_c|\in\Sp_{u^-}$ then
$|P|\in\Sp_{u^-}$ and $|P|\ge q\,p^{u^-}+q$ by (A5). Either way $|P|>q\,p^{e+k'}$; and $q\,p^e$ divides both
$|P|$ (Lemma~\ref{lem:div}) and $q\,p^{e+k'}$; so $|A|=|P|\ge q\,p^{e+k'}+q\,p^e$.

\emph{Case ($\beta$).} Here $|E^-|=e-1$ and $u^-=e-1+k'$. $P$ has two children $N_0,N_1$ of one total $s$,
$|P|=2s$, and $q$ is odd, so $q\mid s$. Their vectors sum to a constant, so if one were balanced both would
be, with equal shapes; hence both are unbalanced. $N_c$ is one of them, it has a witness at its level
$\gamma\in F^-$, and no level $l<\gamma$ lies in $E^-\setminus F^-$. By Lemma~\ref{lem:Spp}, $s\in\Sp_{u^-}$, so
$s\ge q\,p^{u^-}+q$ by (A5). As $q\,p^{e-1}$ divides $s$ and $q\,p^{u^-}$,
$s\ge q\,p^{u^-}+q\,p^{e-1}$ and $|A|=2s\ge q\,p^{e+k'}+q\,p^e$.
\end{proof}

The clauses enter as follows. (2): Step~1 needs that no level of $E$ is less than $a_0$; if there is one, a second nonempty class of
depth $a_0-1$ is forced only through the factor $p$ of that level, which is already in $q\,p^{e+k'}$.
(3): in case ($\beta$) with $p\ge3$ two other children can carry the unequal shapes, so the child $N_c$
that contains the witness may be balanced, and then nothing forces $s$ above $q\,p^{u^-}$. (4):
Lemma~\ref{lem:Spp} needs its hypothesis.

\section{Constructions}\label{sec:constructions}

Elements of $\Z_{p^n}$ are written as digit strings. For a set $T$ of levels let $D(T)$ be the set of
strings whose digits are $0$ outside $T$; for $b\notin T$ and a digit $c$ let $D(T;b\mapsto c)$ be the set of
strings with $u_b=c$ and $0$ at the other levels outside $T$. For a set $D$ of strings and $m\ge1$,
$m\cdot\mathrm{Fib}(D)$ is the multiset with $m$ copies of $(u,w)$ for every $u\in D$ and every layer $w$;
every class has a constant vector in it. $e_t$ is the unit vector of layer $t$.

\subsection{The unbalanced half} Given $E,F\subseteq[1,n]$, let $\tau(u)=1$ if $E\setminus F\ne\emptyset$
and every digit of $u$ at a level of $E\setminus F$ is nonzero, and $\tau(u)=0$ otherwise. Put
\[ A_1=\{(u,\tau(u)):u\in D(E\cup F)\},\qquad A_2=\mathrm{Fib}(D(E)), \]
and $A=A_1$ if $F\subseteq E$, $A=A_1+A_2$ otherwise. Then $|A|$ is $p^{e}$, respectively $p^{e+k}+q\,p^{e}$.

\begin{proposition}\label{prop:CU}
$A$ is unbalanced, $E(A)=E$ and $F(A)=F$.
\end{proposition}
\begin{proof}
Every class that meets $A$ meets $A_1$; its $A_1$-part has a power of $p$ elements, which $q$ does not
divide, and its $A_2$-part is a constant vector; so every nonempty class is unbalanced. At a level of $F$ the
$A_1$-parts of the children of a class have the same layer vector (the digit at that level does not enter
$\tau$) and the $A_2$-parts are constant: equal shapes. At a level of $E$ both parts split evenly: equal
totals. At a level of $F\setminus E$, in the class with all earlier digits $0$, the child with digit $0$
holds the whole $A_2$-part and the others none, while the $A_1$-parts are equal: unequal totals. At a level
outside $E\cup F$ a nonempty class has one nonempty child, and it is unbalanced: unequal totals and unequal
shapes. At a level $a\in E\setminus F$ take a class of depth $a-1$ that meets $A$ and whose digits at the
earlier levels of $E\setminus F$ are nonzero: the $A_1$-part of its child with digit $0$ lies in layer $0$,
that of the child with digit $1$ has the same size and reaches layer $1$, and the $A_2$-parts are constant;
so the two do not have equal shapes.
\end{proof}

\subsection{The balanced half} If $F=[1,n]$, $A=\mathrm{Fib}(D(E))$ has $q\,p^e$ elements, every class
balanced, equal totals exactly at the levels of $E$.

Let $F\ne[1,n]$, with $a_0$, $G$, $k'$ as in the theorem. The building block is a \emph{shape part}
$S=S(a_0,T_1,T_2,H,\text{mode})$, for $T_1\subseteq[1,a_0-1]$, $H\subseteq T_2\subseteq[a_0+1,n]$ and a mode
$\mathrm U$, $\mathrm Q$ or $\mathrm W$. Let $D^*$ be the strings that are $0$ outside
$T_1\cup\{a_0\}\cup T_2$ (the digit of level $a_0$ is free); $\tau(u)=1$ if $H\ne\emptyset$ and every digit
of $u$ at a level of $H$ is nonzero, and $0$ otherwise; and, with $(\Delta x)(w)=x(w)-x(w-1)$,
\begin{align*}
\text{mode U:}\quad&\Lambda_0x=x,\qquad \Lambda_1x=|x|\one-x;\\
\text{modes Q, W:}\quad&\Lambda_0x=|x|\one+\Delta x,\qquad\Lambda_1x=|x|\one-\Delta x;\\
\text{mode Q also:}\quad&\Lambda_ix=|x|\one\quad(2\le i\le p-1).
\end{align*}
$S$ has, for every $u\in D^*$ whose digit $i=u_{a_0}$ has a $\Lambda_i$ defined, and every layer $w$,
$(\Lambda_ie_{\tau(u)})(w)$ copies of $(u,w)$. (All these numbers are $\ge0$: the vectors are $e_t$,
$\one-e_t$, $\one+e_t-e_{t+1}$, $\one-e_t+e_{t+1}$, $\one$.) By linearity the vector of a class of depth $\ge a_0$ with digit $i$ at level $a_0$ is $\Lambda_i\beta$, where $\beta$ counts the strings of $D^*$ in
the class with $\tau=0$ and with $\tau=1$ in layers $0$ and $1$; $\beta$ is supported on layers $0$ and $1$, and for a class that meets $D^*$ its total is a power of
$p$ (odd, if $q=2$), so $\beta$ is not constant. If $\Delta x$ is constant it is $0$ (its entries sum
to $0$) and $x$ is constant; hence, for $i\in\{0,1\}$, $\Lambda_ix-\Lambda_iy$ is constant only if
$x-y$ is, and $\Lambda_0x-\Lambda_1x$, which is $2x-|x|\one$ in mode U and $2\Delta x$ in modes Q, W,
is constant only if $x$ is. For every $x$, $\sum_i\Lambda_ix$ is $|x|\one$, $p|x|\one$, $2|x|\one$ in
modes U, Q, W. Shapes differ in three places. At $a_0$: the children with digits $0,1$ of the class
with all earlier digits $0$ have vectors $\Lambda_0\beta$, $\Lambda_1\beta$ with one $\beta$. At
$a\in H$: in the class of depth $a-1$ with digit $1$ at the levels of $H$ before $a$ and $0$ elsewhere,
the child with digit $0$ has $\beta=p^he_0$ and the child with digit $1$ has $\beta=(p^h-t)e_0+te_1$,
$t\ge1$, $h=|T_2\cap(a,n]|$. At $a>a_0$ outside $T_2$: the class with all earlier digits $0$ has one
child with a non-constant vector and the others empty. With this one reads off (details in Lemmas 9
and 10 of the working note's proof):

\begin{lemma}\label{lem:shape}
Let $\mu=1,p,2$ in modes $\mathrm U,\mathrm Q,\mathrm W$. Then $|S|=\mu\,q\,p^{|T_1|+|T_2|}$; every class of
depth $<a_0$ is balanced in $S$; and for every sum $\Psi$ of multisets $m\cdot\mathrm{Fib}(D)$ the multiset
$A=S+\Psi$ is balanced with $F(A)=[1,a_0-1]\cup(T_2\setminus H)$. In $S$: at a level of $T_1$ and at a level
of $T_2\setminus H$ the children of a class have equal vectors; at a level of $H$ they have equal totals;
at a level $a\ne a_0$ outside $T_1\cup T_2$ only the child with digit $0$ can be nonempty; at level $a_0$ the
children with digits $0,1$ of a class that meets $D^*$ have totals $p^{|T_2|}$ and $(q-1)p^{|T_2|}$ in mode
$\mathrm U$ and $q\,p^{|T_2|}$ each in modes $\mathrm Q$, $\mathrm W$, the remaining children having total
$q\,p^{|T_2|}$ in mode $\mathrm Q$ and $0$ in modes $\mathrm U$, $\mathrm W$.
\end{lemma}

Put $T_2^0=(E\cup F)\cap(a_0,n]$, $H^0=(E\setminus F)\cap(a_0,n]$, and mode$^0=\mathrm Q$ if $a_0\in E$,
$\mathrm U$ if not; so $T_2^0\setminus H^0=F\cap(a_0,n]$. Write $E_<=E\cap[1,a_0-1]$. Define $A$ by the first
case that applies:
\begin{itemize}
\item[(I)] $k'=0$: $A=S(a_0,E_<,T_2^0,H^0,\text{mode}^0)$.
\item[(III)] $k'\ge1$ and $E_<\ne\emptyset$; let $b=\min E$:
 \[A=S(a_0,E_<\setminus\{b\},T_2^0,H^0,\text{mode}^0)+\sum_{i=1}^{p-1}p^{k'}\cdot\mathrm{Fib}(D(E\setminus\{b\};b\mapsto i)).\]
\item[(IV)] $k'\ge1$, $a_0\in E$, $p\ge3$:
 \[A=S(a_0,E_<,T_2^0,H^0,\mathrm W)+\sum_{i=2}^{p-1}p^{k'}\cdot\mathrm{Fib}(D(E\setminus\{a_0\};a_0\mapsto i)).\]
\item[(V)] $k'\ge1$, $a_0\in E$, and some level of $E\setminus F$ lies strictly between $a_0$ and $\min G$;
 let $b$ be the least such:
 \[A=S(a_0,E_<,T_2^0\setminus\{b\},H^0\setminus\{b\},\mathrm Q)+\sum_{i=1}^{p-1}p^{k'}\cdot\mathrm{Fib}(D(E\setminus\{b\};b\mapsto i)).\]
\item[(II)] otherwise: $A=S(a_0,E_<,T_2^0,H^0,\text{mode}^0)+\mathrm{Fib}(D(E))$.
\end{itemize}
In (III), (IV), (V) write $\Psi_i$ for the summand with index $i$. Cases (I), (III), (IV), (V) are clauses (1) to (4), and (II) is reached exactly when $x=1$; the labels are those of the working note and of the Lean files. (The order only fixes which construction is used: each of (III), (IV), (V) gives the right pattern and size wherever its own condition holds.)

\begin{proposition}\label{prop:CB}
In each case $A$ is balanced, $E(A)=E$, $F(A)=F$, and $|A|=q\,p^{e}(p^{k'}+x)$.
\end{proposition}
\begin{proof}[Outline of proof]
The verification level by level is written out in the working note (Section~6 of
\texttt{balanced-half/proof.md}) and is checked in Lean; here is its plan. $F(A)=F$ and balance are Lemma~\ref{lem:shape}, since in every case $T_2\setminus H=F\cap(a_0,n]$. Sizes: in
(I) and (II) the shape part has $q\,p^{e}$, resp.\ $q\,p^{e+k'}$, elements ($\mu=p$ accounts for $a_0\in E$),
and (II) adds $q\,p^e$. In (III), (IV), (V) one level $h$ of $E$ (the \emph{host level}: $b$, $a_0$, $b$) is
treated as if it were not in $E$ when the shape part is built, which then has $q\,p^{e-1+k'}$ elements in
(III), (V) and twice that in (IV); each of the $p-1$ (in (IV): $p-2$) fibre multisets $\Psi_i$ also has $q\,p^{e-1+k'}$ elements, so
$|A|=q\,p^{e+k'}$. At the host level the totals are equal: $S$ and each $\Psi_i$ meet the same classes
of depth $h-1$, those whose digits are $0$ outside $E$, and in each of these every child has
$q\,p^{|E\cap(h,n]|+k'}$ elements, of $S$ in the child with digit $0$ (in (IV): digits $0$ and $1$) and
of $\Psi_i$ in the child with digit $i$. (In (V) this uses that no level of $F\setminus E$ lies between
$a_0$ and $b$, so that $T_2\cap(a_0,b)=E\cap(a_0,b)$ and $|T_2\cap(b,n]|=|E\cap(b,n]|+k'$.) A class of depth $\ge h$ contains either only elements of $S$ or only fibres of
one $\Psi_i$. For $E(A)$: at a level of $E$ other than $h$ both kinds of class split evenly
(Lemma~\ref{lem:shape}; fibres over $D(E\setminus\{h\};\dots)$ are free at the levels of $E\setminus\{h\}$).
At a level $a\ne a_0$ that is outside $E\cup F$, and at a level $a<a_0$ that is not in $E$, the class
with all earlier digits $0$ is nonempty and has one nonempty child. At $a_0\notin E$ (mode U) the children with digits $0$, $1$ of the class with all earlier digits $0$ have totals
$p^{|T_2|}$ and $(q-1)p^{|T_2|}$ in cases (I), (III), and $p^{c}(p^{k'}+q)$ and $p^{c}(q-1)p^{k'}$ with
$c=|E\cap(a_0,n]|$ in case (II), where the fibres lie in the first; any further child is empty. If $p\ge3$ a
nonempty child stands beside an empty one. If $p=2$ then $q$ is odd: in cases (I), (III) $q-1\ne1$, and in case
(II) one total is an odd multiple of $p^c$ and the other an even one.
At a level of $G$: in case (II) the class with all earlier digits $0$ has children with equal $S$-vectors
and all its fibres in the child with digit $0$; in cases (III), (IV), (V) a fibre-only class (digit $1$, $2$,
$1$ at the host level and $0$ elsewhere) is nonempty with all its elements in one child. In case (V) the
level $b$ itself is not in $F(A)$ by Lemma~\ref{lem:shape}, being outside $T_2$.
\end{proof}

Theorem~\ref{thm:main} is Propositions~\ref{prop:U}, \ref{prop:B}, \ref{prop:CU}, \ref{prop:CB} and the first paragraph of Section~4.2, with Lemma~\ref{lem:translate}.

\begin{remark}[one condition instead of four]
For $k'\ge1$, $x=0$ iff the pattern has a \emph{host level}: a level $b\in E$ with $b<\min G$ such that
$b<a_0$, or $b=a_0$ and $p\ge3$, or $b>a_0$, $b\notin F$ and $a_0\in E$. Equivalently (and without a list of
positions): $x=0$ iff some least multiset for the pattern $(E\cup G,F)$, which has $q\,p^{e+k'}$ elements,
contains a nonempty class of depth $<\min G$ all of whose subclasses, itself included, are balanced. Concentrating the fibres
of such a class then makes the totals unequal at the levels of $G$ at no cost. The word ``least'' cannot be dropped: for $p=2$, $q=3$, $n=3$, $E=\emptyset$, $F=\{3\}$, where $x=1$, a multiset of $12$ elements with pattern $(\{3\},\{3\})$ has such a class; the least ones, of $6$ elements, have none. (Propositions~11 and~13 of the working note's proof; the example is its referee's.)
\end{remark}

\section{Larger groups}\label{sec:lift}

Theorem~\ref{thm:main} holds as it stands on $\Z_{p^nq\,s}$ for every $s\ge1$, for the patterns that consist of
divisors of $p^nq$. This rests on two lemmas that are not the author's. The first, on lifting along an odd prime, was found, with another proof, by a fresh instance of the model that
was reading a sequel to this paper; the proof below, by counting roots, is the one that was then formalised (Section~\ref{sec:lean}).
The second, on signs, which gives the prime~$2$, was proved by a further instance (see below), and has since been formalised too.
A short search of the literature, made by the author after the results existed (seven requests), found neither lemma and not
Corollary~\ref{cor:lift}; that is weak evidence, and statements as simple as these may well be in print. The first step of the
proof of Lemma~\ref{lem:lift}, that reduction keeps the divisors, is in \cite[Section~2.2]{KLaMS}.

\begin{lemma}[lifting]\label{lem:lift}
Let $N\ge1$, let $r$ be an odd prime, $M=Nr$, and let $A$ be a nonempty multiset on $\Z_N$. There is a
multiset $A'$ on $\Z_M$ with the same number of elements, whose reduction modulo $N$ is $A$, such that for
every divisor $d$ of $M$: $\Phi_d\mid A'(X)$ if and only if $d\mid N$ and $\Phi_d\mid A(X)$.
\end{lemma}
\begin{proof}
If $A'$ reduces to $A$ modulo $N$, then $A'(X)\equiv A(X)$ modulo $X^N-1$, and $\Phi_d\mid X^N-1$ for $d\mid N$; so
$A'$ and $A$ have the same divisors among the $\Phi_d$, $d\mid N$, whatever the lift. It remains to choose the
lift so that no $\Phi_d$ with $d\mid M$, $d\nmid N$ divides $A'(X)$.

Write $N=r^kN'$ with $r\nmid N'$, and for $y\in\Z_M$ let $\beta(y)\in[0,r^{k+1})$ be its residue modulo
$r^{k+1}$. An element $x$ of $\Z_N$ has $r$ lifts to $\Z_M$, and their values of $\beta$ are
$(x \bmod r^k)+r^kt$, $t=0,\dots,r-1$, each once, because $N'$ is invertible modulo $r$. Lift every element of
$A$ to its lift with $t=0$, except one single element $x_0$, which is lifted to its lift $y_0$ with $t=1$.

A divisor $d$ of $M$ that does not divide $N$ is $d=e\,r^{k+1}$ with $e\mid N'$. Suppose $\Phi_d\mid A'(X)$, and fix a root of
unity $\omega$ of order $e$. For every root of unity $\eta$ of order $r^{k+1}$ the product $\omega\eta$ is a
primitive $d$-th root of unity, so
\[
0=A'(\omega\eta)=\sum_{y}\omega^{y}\eta^{\beta(y)}=\ell(\eta),\qquad
\ell(Y)=\sum_{b}g_b\,Y^{b},\quad g_b=\sum_{\beta(y)=b}\omega^{y},
\]
the sums over the elements $y$ of $A'$ with multiplicity. Every element has $\beta(y)<2r^k$, so $\ell$ has
degree less than $2r^k\le(r-1)r^k$, the number of roots of unity of order $r^{k+1}$, at all of which it
vanishes; hence $\ell=0$. But for $b=\beta(y_0)$, which lies in $[r^k,2r^k)$, the only element with
$\beta(y)=b$ is $y_0$, so $g_b=\omega^{y_0}\ne0$.
\end{proof}

The inequality $2r^k\le(r-1)r^k$ is the only place where the oddness of $r$ is used. Lifting everything to the least
representative and one element higher by $N$ does not always do: for $N=2$, $r=3$ and $A=\{0,0,1\}$, raising one of the two $0$s gives $\{0,1,2\}$ in $\Z_6$,
whose mask is $\Phi_3$.

\subsection*{The prime 2}

For $r=2$ the count of roots fails, since $2\cdot2^k$ is more than the number $2^k$ of roots of unity of order $2^{k+1}$, and the
statement turns into one about signs. The author put that statement as a question to a further fresh instance, which
proved it for the moduli that are needed here (and for some others, see below). Lemma~\ref{lem:signs} is that
instance's, and so are the two steps of its proof, the descent modulo $q$ and the choice of signs; the proof below
reads the coefficients directly where the instance uses a lemma on functions that sum to zero over cosets, and that
part of the wording is the author's. In the two lemmas that follow, $n$ is a modulus and $p$, $q$ are odd primes; they
are not the $n$, $p$, $q$ of Theorem~\ref{thm:main}.

\begin{lemma}[signs]\label{lem:signs}
Let $p\ne q$ be odd primes, $a,b\ge0$, $n=p^aq^b$, and let $B$ be a nonempty multiset on $\Z_n$. Each element of $B$
can be given a sign $\pm1$ (equal elements may get different signs) so that the signed mask
$S(X)=\sum_{x\in B}\pm X^x$ is divisible by no $\Phi_d$ with $d\mid n$, the divisor $d=1$ included.
\end{lemma}
\begin{proof}
For a set $C$ of residues write $s(C)$ for the sum of the signs of the elements of $B$ that lie in $C$.

\emph{What has to be avoided.} Let $d=p^iq^j$ divide $n$ and suppose $\Phi_d\mid S$. If $j\ge1$ then
$\Phi_{p^iq^j}(X)\,\Phi_{p^i}(X^{q^{j-1}})=\Phi_{p^i}(X^{q^j})$, and $f(X^q)\equiv f(X)^q$ modulo $q$ for every integer
polynomial $f$, so $\Phi_{p^iq^j}\equiv\Phi_{p^i}^{\,q^j-q^{j-1}}$ modulo $q$. In every case, therefore, $\Phi_{p^i}$
divides $S$ in $\mathbb{F}_q[X]$, and with it the remainder of $S$ modulo $X^{p^i}-1$, which is
$\sum_{v=0}^{p^i-1}s(v+p^i\Z)\,X^v$. For $i=0$ this says that $s(\Z_n)\equiv0$ modulo $q$. For $i\ge1$ we have
$\Phi_{p^i}=1+X^{p^{i-1}}+\dots+X^{(p-1)p^{i-1}}$; a multiple $g\,\Phi_{p^i}$ of degree less than $p^i$ has
$\deg g<p^{i-1}$, so its coefficient at $X^v$ is the coefficient of $g$ at $X^{v\bmod p^{i-1}}$. So the $p$ classes
modulo $p^i$ that make up a class modulo $p^{i-1}$ all have the same $s$ modulo $q$. It is therefore enough to
choose the signs so that
\begin{itemize}
\item[(i)] $s(\Z_n)\not\equiv0$ modulo $q$, and
\item[(ii)] for each $i=1,\dots,a$ there are two classes modulo $p^i$, inside one class modulo $p^{i-1}$, whose
values of $s$ are incongruent modulo $q$.
\end{itemize}

\emph{The signs.} A class that holds $\nu$ elements of $B$ can be given any value $s$ with $|s|\le\nu$ and
$s\equiv\nu$ modulo $2$. Fix an element $x^*$ of $B$ and let $C_i$ be its class modulo $p^i$, so that
$\Z_n=C_0\supseteq C_1\supseteq\dots\supseteq C_a$. Give $C_a$ the value $s(C_a)=1$ if it holds an odd number of
elements and $2$ if an even number. Then, for $i=a,a-1,\dots,1$ in turn, sign the rest of $C_{i-1}$, which
consists of $p-1$ classes modulo $p^i$. Those that hold an even number of elements (none, perhaps) get $s=0$. Those that hold an
odd number get $s=+1$ and $s=-1$ in pairs; if one is left over, it gets that one of $+1$, $-1$ which keeps
$s(C_{i-1})$ away from $0$ modulo $q$ (if both $s(C_i)+1$ and $s(C_i)-1$ were divisible by $q$, so would $2$ be).
Thus $s(C_{i-1})$ is $s(C_i)$ or $s(C_i)\pm1$, and, starting from $s(C_a)\in\{1,2\}$, no $s(C_i)$ is divisible by
$q$. For $i=0$ this is (i). For (ii) at the level $i$: if one of the $p-1$ classes holds an even number of
elements, its $s$ is $0\not\equiv s(C_i)$. Otherwise all $p-1$ hold odd numbers; $p-1$ is even and at least
$2$, so the values $+1$ and $-1$ both occur among them, and $s(C_i)$ is incongruent to one of the two.
\end{proof}

\begin{lemma}[lifting along 2]\label{lem:lift2}
Let $N=2^kN'$ with $k\ge0$ and $N'=p^aq^b$ as in Lemma~\ref{lem:signs}, and let $A$ be a nonempty multiset on $\Z_N$.
There is a multiset $A'$ on $\Z_{2N}$ with the same number of elements, whose reduction modulo $N$ is $A$, such
that for every divisor $d$ of $2N$: $\Phi_d\mid A'(X)$ if and only if $d\mid N$ and $\Phi_d\mid A(X)$.
\end{lemma}
\begin{proof}
As in Lemma~\ref{lem:lift} the divisors $d\mid N$ take care of themselves, and the others are $d=2^{k+1}e$ with
$e\mid N'$. The two lifts of an element $x$ of $\Z_N$ are congruent modulo $2^{k+1}$ to $x\bmod 2^k$ (the
\emph{low} lift) and to $(x\bmod2^k)+2^k$ (the \emph{high} one), because $N'$ is odd. Choose a class $c$ modulo
$2^k$ that holds an element of $A$, let $B$ be the multiset on $\Z_{N'}$ of the residues modulo $N'$ of the
elements of $A$ in that class, and give them signs by Lemma~\ref{lem:signs}. Lift the elements with the sign
$+$ low and those with the sign $-$ high; lift the elements of all other classes low.

Suppose $\Phi_d\mid A'(X)$ with $d=2^{k+1}e$, and fix a root of unity $\omega$ of order $e$. For each of the $2^k$
roots of unity $\eta$ of order $2^{k+1}$ the product $\omega\eta$ is a primitive $d$-th root of unity, and
$\eta^{2^k}=-1$, so
\[
0=A'(\omega\eta)=\sum_{c'=0}^{2^k-1}T_{c'}\,\eta^{c'},\qquad T_{c'}=\sum_{y\equiv c'\ (2^k)}\pm\,\omega^{y},
\]
the inner sum over the elements $y$ of $A'$ in the class $c'$ modulo $2^k$, with the sign $-$ for the high lifts.
The polynomial $\sum_{c'}T_{c'}Y^{c'}$ has degree less than $2^k$ and vanishes at $2^k$ points, so every
$T_{c'}$ is $0$. But $T_c=S(\omega)$, where $S$ is the signed mask of $B$, and $\omega^{N'}=1$; this contradicts
Lemma~\ref{lem:signs}.
\end{proof}

For odd $N$ the statement of Lemma~\ref{lem:lift2} and that of Lemma~\ref{lem:signs} for $n=N$ are the same
statement. The proof of Lemma~\ref{lem:lift2} uses nothing of $N'$ but Lemma~\ref{lem:signs} for $n=N'$; so
lifting along $2$ holds for every $N$ for whose odd part the statement on signs holds. The same instance proved
that statement also for $n=p\,r\,q^b$ with three different odd primes, starting from a signing of a rectangle of cells
whose row sums and column sums are all $0$ or $\pm1$, in which one sign is then turned if need be; that proof is not
given here and is not formalised (it is in the file \texttt{sign-lemma-result.md} that goes with the second working
note, and the author has read it). Together the
two cases cover every odd $n$ below $1155=3\cdot5\cdot7\cdot11$.

The author then put the general statement to a further fresh instance. It proved it for every squarefree odd $n$,
and for every odd $n$ in which at most one prime is repeated, and for some more (its Theorem~6, in the file
\texttt{sign-general-result.md} that goes with the second working note); with the two cases above this covers every
odd $n$ below $118125=3^3\cdot5^4\cdot7$. The signs there are of product form: for squarefree $n$ and a set $B$,
the sign of $x$ is $\prod_{p\mid n}\sigma_p(x\bmod p)$ with functions $\sigma_p\colon\Z_p\to\{\pm1\}$, and the proof
counts, by an induction that splits off one prime at a time, that at least the fraction $2^{-k}$ of all such choices
is good, $k$ being the number of primes. The author has read the proofs for squarefree $n$ and for one repeated prime and found no error, and has not checked the further cases, on which the bound $118125$ rests from $n=1575=3^2\cdot5^2\cdot7$ on. By that instance's account another instance, which it asked, read the proofs too; and a further fresh instance, asked by the author to referee them, read all of them line by line, the further cases included, tested them with programs of its own, recomputed the bound, and found no error. None of it is formalised. For general odd $n$ the statement on signs is
open, and with it Lemma~\ref{lem:lift} for $r=2$ and general $N$; no counterexample was found. A stronger statement
is false: in every search of the author's it had been enough to take all signs $+$ or to turn one sign, and the instance that proved the above gives a multiset on $\Z_n$, $n=3\cdot5\cdot7\cdot11\cdot13\cdot17\cdot19\cdot23\cdot29\cdot103$, for which the mask
with all signs $+$ is divisible by $\Phi_{103}$ and every mask with one sign turned by $\Phi_p$ for a prime $p\le29$.

\subsection*{The theorem on larger groups}

\begin{corollary}\label{cor:lift}
Let $p\ne q$ be primes, $n\ge0$, and $s\ge1$; put $M=p^nq$. Let $D$ be a set of divisors $d>1$ of $M$. The numbers of
elements of the nonempty multisets on $\Z_{Ms}$ whose mask is divisible by $\Phi_d$, among the divisors $d>1$ of
$Ms$, exactly for $d\in D$, are the numbers of elements of the nonempty multisets on $\Z_M$ with pattern exactly
$D$. In particular the least of them is the number that Theorem~\ref{thm:main} gives for $D$.
\end{corollary}
\begin{proof}
For $s$ a prime $r$: a multiset on $\Z_{Mr}$ as described reduces modulo $M$ to a multiset of the same
size on $\Z_M$ with pattern exactly $D$; and a nonempty multiset on $\Z_M$ with pattern $D$ has, by
Lemma~\ref{lem:lift} or, for $r=2$, by Lemma~\ref{lem:lift2}, a lift of the same size whose divisors among the $\Phi_d$, $d\mid Mr$, are
those $d\in D$. For general $s=2^js'$ with $s'$ odd, apply this once for each prime factor of $s$, counted with
multiplicity, and take the factors $2$ first: the groups lifted from are then $\Z_{2^iM}$, $i<j$, whose order has the
odd part of $M$, which is of the form $p^aq^b$ of Lemma~\ref{lem:signs} (a prime power if $p$ or $q$ is $2$), so that
Lemma~\ref{lem:lift2} applies; the odd prime factors follow by Lemma~\ref{lem:lift}, which holds for every $N$.
\end{proof}

The lower bound in the corollary needs nothing about the new divisors: a nonempty multiset on $\Z_{Ms}$ whose mask is
divisible by $\Phi_d$, among the divisors $d>1$ of $M$, exactly for $d\in D$, reduces modulo $M$ to one of the same
size with pattern exactly $D$, whatever $\Phi_d$ with $d\nmid M$ divide it; so it has at least as many elements as
Theorem~\ref{thm:main} says.

For $r=p$ the last sentence of the corollary says nothing new: $\Z_{p^nq\cdot p}=\Z_{p^{n+1}q}$, and it is Theorem~\ref{thm:main} for $n+1$ and the patterns
without the level $n+1$ (this was pointed out by Noon, who read the formal statements). The new ground is $r=q$ and a third prime.
For instance, on $\Z_{p^2q^2}$ the $32$ patterns among the divisors of $p^2q$ have the least sizes of
Theorem~\ref{thm:main} for $n=2$, and those among the divisors of $pq^2$ the same with $p$ and $q$ exchanged.

\subsection*{Sets on larger groups}

Theorem~\ref{thm:main} is about multisets, and on $\Z_{p^nq}$ itself a pattern may have no set at all
(Section~\ref{sec:intro-sets}). On a larger group the copies of an element can be sent to different lifts, and a
least multiset can become a set. Lifting a multiset on $\Z_N$ to a set on a larger group, so that the divisors
$\Phi_s$, $s\mid N$, are kept, is in \cite{KLaMS}: in its Section~2.1 the copies of $a$ go to $a,a+N,a+2N,\dots$,
with no attention to the other divisors (that lift can gain one, as in the example after Lemma~\ref{lem:lift}),
and in its Section~8 the lift is chosen to gain a divisor; the point here is a lift that gains none.

\begin{proposition}[lifting to a set]\label{prop:setlift}
Let $r$ be an odd prime, $N\ge1$, $j\ge1$, and let $A$ be a nonempty multiset on $\Z_N$ in which every element has
multiplicity at most $r^j$ and not every element has multiplicity $r^j$. There is a set $A'\subseteq\Z_{Nr^j}$
with as many elements as $A$, whose reduction modulo $N$ is $A$, such that for every divisor $d$ of $Nr^j$:
$\Phi_d\mid A'(X)$ if and only if $d\mid N$ and $\Phi_d\mid A(X)$.
\end{proposition}
The condition cannot be dropped: if every element has multiplicity $r^j$, the only set on $\Z_{Nr^j}$ that reduces to $A$ is
the full preimage of its support, whose mask is divisible by every $\Phi_d$ with $d\mid Nr^j$, $d\nmid N$.
\begin{proof}
\emph{One step.} Let $i\ge1$, and let $A$ be a nonempty multiset on $\Z_N$ in which every element has multiplicity
at most $r^i$, and not every element multiplicity $r^i$. We lift it to a multiset $A_1$ on $\Z_{Nr}$ with the same number of elements and no new
divisor, in which every element has multiplicity at most $r^{i-1}$ and, if $i\ge2$, not every element multiplicity $r^{i-1}$. Doing this for
$i=j,j-1,\dots,1$, on $\Z_N,\Z_{Nr},\dots$, ends with multiplicities at most $1$.

Keep the notation of the proof of Lemma~\ref{lem:lift}: $N=r^kN'$, and the $r$ lifts of $x$ have
$\beta=(x\bmod r^k)+r^kt$, $t=0,\dots,r-1$; call $t$ the \emph{slot} of the lift. A lift of $A$ is given by
numbers $a_x(t)\ge0$ with $\sum_ta_x(t)=m(x)$, the multiplicity of $x$: $a_x(t)$ copies of $x$ go to its lift in
the slot $t$. Let $d=e\,r^{k+1}$ with $e\mid N'$, let $\omega$ have order $e$, and suppose $\Phi_d\mid A_1(X)$.
As in that proof $\ell(\eta)=0$ for every root of unity $\eta$ of order $r^{k+1}$, where now $\ell$ may have any
degree below $r^{k+1}$. So $\ell$ is divisible by $\prod_\eta(Y-\eta)=1+Y^{r^k}+\dots+Y^{(r-1)r^k}$, with a quotient
of degree less than $r^k$, and therefore $g_{c+r^kt}$ does not depend on $t$ ($g_{c+r^kt}=G_c(t)$ below, because
$\omega^y=\omega^x$ when $y$ lifts $x$): for every class $c$ modulo $r^k$,
\[ G_c(t)=\sum_{x\equiv c\ (r^k)}a_x(t)\,\omega^{x}\quad\text{is the same for all slots } t. \]
Hence there is no new divisor as soon as, for one element $x_1$ and its class $c$,
\[ a_{x_1}(0)-a_{x_1}(r-1)\in\{1,2\}\quad\text{and}\quad a_x(0)=a_x(r-1)\ \text{ for every other } x\equiv c\ (r^k), \]
since then $G_c(0)-G_c(r-1)$ is $\omega^{x_1}$ or $2\omega^{x_1}$, for every $\omega$.

Write $m(x)=ra+\rho$ with $0\le\rho<r$, and call the slots $1,\dots,r-2$ the middle slots; there is at least one, because
$r\ge3$. For $x\ne x_1$ let $a_x(t)=a+1$ on $\rho$ slots and $a$ on the others, the $\rho$ slots being middle slots
if $\rho\le r-2$, and all slots but the slot $1$ if $\rho=r-1$; in both cases $a_x(0)=a_x(r-1)$. If some
multiplicity is not divisible by $r$, let $x_1$ be such an element, and put $a_{x_1}(0)=a+1$, $a_{x_1}(r-1)=a$,
and $a+1$ on $\rho-1$ middle slots, $a$ on the other middle slots. Otherwise let $x_1$ be an element with
$m(x_1)=ra<r^i$, and put $a_{x_1}(0)=a+1$, $a_{x_1}(r-1)=a-1$, and $a$ on the middle slots.

The new multiplicities are at most $\lceil m(x)/r\rceil\le r^{i-1}$, and in the last case $a+1\le r^{i-1}$ because
$ra<r^i$. If $i\ge2$, one lift of $x_1$ has a multiplicity that is positive and less than $r^{i-1}$: in the first
case the one in the slot $0$ if $a+1<r^{i-1}$, and the one in the slot $r-1$ if $a+1=r^{i-1}$ (then $a\ge1$); in
the last case one in a middle slot (there $a\ge1$).
\end{proof}

\begin{corollary}[sets]\label{cor:sets}
Let $p\ne q$ be primes, $n\ge0$, $M=p^nq$, let $D$ be a set of divisors $d>1$ of $M$, and let $m$ be the number
that Theorem~\ref{thm:main} gives for $D$. If $s$ is divisible by $r^j$ for an odd prime $r$ with $r^j>m$, then
the least number of elements of a \emph{set} in $\Z_{Ms}$ whose mask is divisible by $\Phi_d$, among the divisors
$d>1$ of $Ms$, exactly for $d\in D$, is $m$.
\end{corollary}
\begin{proof}
A set is a multiset, so it has at least $m$ elements by Corollary~\ref{cor:lift}. For the other direction take a
multiset of $m$ elements on $\Z_M$ with pattern exactly $D$; its multiplicities are at most $m<r^j$. Lift it along
the factors $2$ of $s$ as in the proof of Corollary~\ref{cor:lift} (a lift does not raise a multiplicity), then along
$r^j$ by Proposition~\ref{prop:setlift}, which makes it a set, then along the remaining prime factors of $s$, counted
with multiplicity (they are odd; further factors $r$ may be among them), by Lemma~\ref{lem:lift}; a lift of a set is a set.
\end{proof}

The bound $r^j>m$ is crude: what is used is that some least multiset has its multiplicities at most $r^j$ (they
cannot all be equal to a number greater than $1$, or dividing them by it would give a smaller multiset with the
same pattern). The least multisets of Section~\ref{sec:constructions} for the unbalanced half have multiplicities at
most $2$, so there every $s$ with an odd prime factor will do. In the ten groups of Section~\ref{sec:intro-sets} a solver
found for each of the $1{,}088$ patterns a least multiset with multiplicities at most $3$ (at most $2$ when
$p$ is odd; the author's run left one pattern of $\Z_{48}$ undecided, and the referee of this subsection found the last with a model of its own); and each of the $96$ least multisets in the solver's tables for
$\Z_{12}$, $\Z_{18}$ and $\Z_{20}$ was lifted by the proof's recipe with $r=3$, and the $64$ of $\Z_{12}$ and $\Z_{18}$ again with $r=5$, every time to a set with the same
pattern, checked by exact division. For instance, exactly $\Phi_6$ among the divisors needs five elements; on $\Z_6$ every such
multiset has a repeated element, and on $\Z_{18}$ the set $\{0,1,2,4,10\}$ will do (here $r^j=3$ is less than $m=5$: the multiplicities, at most $2$, are what matters).

\emph{The prime $2$.} For $r=2$ the proposition is false as it stands: on $\Z_2$ the multiset with $0$ four times and $1$
twice has no cyclotomic divisor, and each of its six lifts to a set on $\Z_8$ gains $\Phi_4$ or $\Phi_8$. One step along $2$ does
work. Let $N=2^kN'$ with $N'$ as in Lemma~\ref{lem:signs}, and let every element of $A$ have multiplicity $1$ or $2$, not every
element $2$. Send the two copies of a doubled element to its two lifts; in the sums $T_{c'}$ of the proof of
Lemma~\ref{lem:lift2} they cancel. Choose a class modulo $2^k$ that holds an element of multiplicity $1$, sign the elements of
multiplicity $1$ in it by Lemma~\ref{lem:signs}, and lift them low or high by their signs; those of the other classes may go to either lift; the proof of Lemma~\ref{lem:lift2}
then shows that the lift, which is a set, has no new divisor. So the conclusion of Corollary~\ref{cor:sets} also holds for every even $s$ when some least multiset has multiplicities at most $2$: not every element of it has multiplicity $2$, as above; make the first lift along $2$ in this way, which gives a set, and lift along the remaining prime factors of $s$ as in the proof of Corollary~\ref{cor:lift}; a lift of a set is a set. In the ten groups that is so for every pattern when $p$ is odd. Without
such a condition there may be no set at all: on $\Z_{24}=\Z_{12\cdot2}$ the balanced pattern $E=\{1\}$, $F=\{2\}$ of $\Z_{12}$,
with least multisets of $18$ elements, has none.

Proposition~\ref{prop:setlift}, Corollary~\ref{cor:sets} and the remarks on the prime $2$ are the author's; they are not
formalised, and were found after the rest of this section had been read. Two requests by the author did not find the
proposition in print (the second working note lists them); that is weak evidence.

\section{The formal proof}\label{sec:lean}

Theorem~\ref{thm:main} is formalised in Lean~4 (toolchain \lean{v4.37.0-rc1}) with Mathlib, in the library
\lean{Pnq} of the development that also holds the tiling theorems of the companion papers. The eleven files of Theorem~\ref{thm:main} are in the shared development (\lean{lean-tiling/Pnq/}), and copies are in
\lean{minsize-agent3/lean-min/}, beside the working note, with a log of what was checked, the statements as they were fixed before any proof that uses them, and the audit of the statements. The author
built the whole development with them, and then again from clean in a fresh copy of the shared folder; Noon then rebuilt it from clean in a copy of its own, with the same result (in both, Mathlib itself was not rebuilt).

\subsection*{The statements} A multiset on $\Z_{p^nq}$ is a function \lean{w : ZMod (p \^{} n * q) → ℕ}, its
multiplicities. Its mask polynomial, with exponents the representatives in $[0,M)$, is
{\footnotesize
\begin{verbatim}
noncomputable def maskW (w : ZMod M → ℕ) : Polynomial ℤ :=
  ∑ a : ZMod M, Polynomial.C (w a : ℤ) * Polynomial.X ^ a.val
\end{verbatim}
}
and \lean{levelsEW w}, \lean{levelsFW w} are the sets of $\ell\in[1,n]$ with
\lean{Polynomial.cyclotomic (p \^{} ℓ) ℤ ∣ maskW w}, respectively
\lean{Polynomial.cyclotomic (p \^{} ℓ * q) ℤ ∣ maskW w}. The formulas of the theorem are four definitions:
{\footnotesize
\begin{verbatim}
def gOf (Es Fs : Finset ℕ) (a0 : ℕ) : Finset ℕ := (Fs \ Es).filter (fun ℓ => a0 < ℓ)

def xOf (p : ℕ) (Es Fs : Finset ℕ) (a0 : ℕ) : ℕ :=
  if gOf Es Fs a0 = ∅ ∨ (∃ ℓ ∈ Es, ℓ < a0) ∨ (a0 ∈ Es ∧ 3 ≤ p) ∨
      (a0 ∈ Es ∧ p = 2 ∧ ∃ b ∈ Es \ Fs, a0 < b ∧ ∀ g ∈ gOf Es Fs a0, b < g)
  then 0 else 1

def minU (p q : ℕ) (Es Fs : Finset ℕ) : ℕ :=
  if Fs ⊆ Es then p ^ Es.card else p ^ Es.card * (p ^ (Fs \ Es).card + q)

def minB (p q : ℕ) (Es Fs : Finset ℕ) (a0 : ℕ) : ℕ :=
  q * p ^ Es.card * (p ^ (gOf Es Fs a0).card + xOf p Es Fs a0)
\end{verbatim}
}
and its three parts are three theorems (\lean{Pnq/MinGroup.lean}):
{\footnotesize
\begin{verbatim}
theorem minsize_unbalanced_zmod (p q n : ℕ) [Fact p.Prime] [Fact q.Prime] (hpq : p ≠ q)
    (Es Fs : Finset ℕ) (hE : Es ⊆ Finset.Icc 1 n) (hF : Fs ⊆ Finset.Icc 1 n) :
    IsLeast {s | ∃ w : ZMod (p ^ n * q) → ℕ, ∑ a, w a = s ∧
        ¬ Polynomial.cyclotomic q ℤ ∣ maskW w ∧ levelsEW w = Es ∧ levelsFW w = Fs}
      (minU p q Es Fs)

theorem minsize_balanced_all_zmod (p q n : ℕ) [Fact p.Prime] [Fact q.Prime] (hpq : p ≠ q)
    (Es : Finset ℕ) (hE : Es ⊆ Finset.Icc 1 n) :
    IsLeast {s | ∃ w : ZMod (p ^ n * q) → ℕ, ∑ a, w a = s ∧ 0 < s ∧
        Polynomial.cyclotomic q ℤ ∣ maskW w ∧ levelsEW w = Es ∧
        levelsFW w = Finset.Icc 1 n}
      (q * p ^ Es.card)

theorem minsize_balanced_zmod (p q n : ℕ) [Fact p.Prime] [Fact q.Prime] (hpq : p ≠ q)
    (Es Fs : Finset ℕ) (hE : Es ⊆ Finset.Icc 1 n) (hF : Fs ⊆ Finset.Icc 1 n)
    (a0 : ℕ) (ha1 : 1 ≤ a0) (han : a0 ≤ n)
    (hlow : ∀ ℓ, 1 ≤ ℓ → ℓ < a0 → ℓ ∈ Fs) (hna : a0 ∉ Fs) :
    IsLeast {s | ∃ w : ZMod (p ^ n * q) → ℕ, ∑ a, w a = s ∧
        Polynomial.cyclotomic q ℤ ∣ maskW w ∧ levelsEW w = Es ∧ levelsFW w = Fs}
      (minB p q Es Fs a0)
\end{verbatim}
}
\noindent (line breaks differ from the files). \lean{IsLeast S m} is Mathlib's ``$m\in S$ and $m\le s$ for
every $s\in S$'', so each theorem has both halves of its part of Theorem~\ref{thm:main}: a multiset with exactly
the pattern and $m$ elements exists, and none with exactly the pattern has fewer. In the third, the
hypotheses on \lean{a0} say that it is the least level in $[1,n]$ outside \lean{Fs}. The three together cover
every pattern once.

\subsection*{The parts} \sloppy On the tree of Section~\ref{sec:tree} (\lean{Pnq/Tree.lean}: a $p$-ary tree of depth
$n$ whose leaves are functions \lean{Fin q → ℕ}, with $E$ and $F$ defined by recursion as in the text), the
lower bounds are \lean{pow\_card\_union\_le\_size}, \lean{sharp\_lemma3} (\lean{Pnq/Sharp.lean}) and
\lean{balanced\_base}, \lean{balanced\_plus} (\lean{Pnq/Balanced.lean}); the constructions are
\lean{exists\_unbalanced} (\lean{Pnq/BuildU.lean}) and \lean{exists\_balanced\_all}, \lean{exists\_balanced}
(\lean{Pnq/BuildB.lean} with five files of lemmas); \lean{Pnq/MinTree.lean} joins them into three
\lean{IsLeast} statements about trees. All of these hold for coprime $p,q\ge2$, prime or not.
\lean{Pnq/Weighted.lean} is the bridge to the group: the Fourier zeros of a multiset, their equivalence with
cyclotomic divisibility of the mask polynomial (\lean{zeroAtW\_iff\_cyclotomic\_dvd}), the tree
\lean{treeOfW w} of a multiset, the translation of Lemma~\ref{lem:translate} (\lean{size\_treeOfW},
\lean{E\_treeOfW}, \lean{F\_treeOfW\_zero}, \lean{F\_treeOfW}), and that every tree is the tree of a
multiset (\lean{treeOfW\_surjective}; the Chinese remainder theorem); it also proves that on a set these
notions are the ones used for the tiling theorems. The eleven files have $3{,}484$ lines.

\subsection*{How the formal proof differs from this paper} The formal lower bound for part (a) is not
Proposition~\ref{prop:U}: it is the direct induction of the working note (two claims and four cases), written
before the argument by sets of sizes was found. The formal lower bounds for part (c) follow
Lemmas~\ref{lem:shapes} to~\ref{lem:Spp}, Lemma~\ref{lem:bal} and Proposition~\ref{prop:B}. The formal constructions are recursive
builders of trees that follow Section~\ref{sec:constructions}, with changes of convenience that the instances
that wrote them reported: a string counts as having all its digits at the levels of $H$ (or of
$E\setminus F$) nonzero also when there are no such levels, which renames layers; in modes $\mathrm Q$, $\mathrm W$, for $q\ge3$, the vector for digit $0$ on the strings with $\tau=1$ is
$\one+e_0-e_2$, where this paper has $\one+e_1-e_2$; and in cases (III) and (V) any level that satisfies the condition serves as host, not only the least. The remark on host levels at the end of
Section~\ref{sec:constructions} is not formalised.

\subsection*{How it was made, and what was checked} The author fixed every definition and every theorem
statement first, in files that compiled with the proofs left open, and before any proof in the new files was written had Lean evaluate \lean{minU}, \lean{minB} and $q\,p^{e}$ on the $1{,}265$ patterns whose least size the solvers had
decided (Section~\ref{sec:computations}): all agree. Two fresh instances of the model had written the proofs of the lower bounds on the tree in this way earlier (\lean{Sharp.lean}, \lean{Balanced.lean}); four more then each wrote the
proofs of one of the remaining parts in a copy of its own, forbidden to change a statement. On the merged copy the author
checked that every fixed declaration stands character for character; that Lean prints each fixed definition
and statement exactly as it printed them before the proofs existed; that the new files contain no
\lean{sorry}, no added axiom, and no instance, notation, macro or attribute that could change what a statement
means; that the whole development builds; and that \lean{\#print axioms} gives, for each of the theorems named
above, only \lean{propext}, \lean{Classical.choice} and \lean{Quot.sound}. Two readings of the \emph{statements} followed. Noon read \lean{Rule.lean} and \lean{MinGroup.lean} whole, and in \lean{Weighted.lean} the definitions
\lean{maskW}, \lean{hatW}, \lean{ZeroAtW} and the statement of \lean{zeroAtW\_iff\_cyclotomic\_dvd},
against the rule of the working note, and found the four clauses to be the note's word for word; it did
not read the constructions or the rest of \lean{Weighted.lean}; it read \lean{MinGroup.lean} before three of its comments were reworded on the audit's advice (no statement changed), and again afterwards. Its own trial, by integer
programming on multisets, gave the rule's value in all $1{,}114$ cases it tried, in thirteen groups ($\Z_6$ taken with both choices of $(p,q)$; the solver works in floating point; the multisets it found were checked in whole numbers). A fresh instance was asked to audit the three statements
and the seven definitions and to try to show that they say something other than intended: it read them in
Lean's printout with every instance shown and found no difference from the claim; proved in Lean that the
hypotheses on \lean{a0} hold for exactly one number when $F\ne[1,n]$ and for none otherwise, so that the three
statements cover every pattern once; compared Lean's values with the rule as worded in the working note on
$103{,}752$ patterns; and, by enumerating multisets with exact polynomial division, found Lean's value to be
the least size on every pattern of $\Z_2,\Z_3,\Z_5,\Z_6,\Z_{10},\Z_{15},\Z_{12},\Z_{18},\Z_{20},\Z_{24}$ and on
every pattern it could settle in nine further groups: $936$ cases, over the $28$ triples $(p,q,n)$ it took in these $19$ groups, no mismatch. It did not read the proofs.

\subsection*{Larger groups} The three lemmas and Corollary~\ref{cor:lift} of Section~\ref{sec:lift} are formalised in six further files. Three hold the odd primes,
\lean{Pnq/Lift.lean}, \lean{Pnq/Lift1.lean} and \lean{Pnq/LiftMin.lean} (603 lines):
\lean{Pnq.Lift.sum\_reduce} and \lean{Pnq.Lift.cyclotomic\_dvd\_reduce\_iff} (reduction modulo $N$ keeps the number of elements and the divisors $\Phi_d$, $d\mid N$);
\lean{Pnq.Lift.exists\_lift} (for $r\neq2$ a nonempty multiset on \lean{ZMod N} is the reduction of one on
\lean{ZMod (N * r)} whose mask is divisible by no $\Phi_d$ with $d\mid Nr$, $d\nmid N$);
\lean{Pnq.Lift.sizes\_eq} (for $r\neq2$ and a set $P$ of divisors of $N$ the sizes of the nonempty multisets with exactly the
pattern $P$ are the same on the two groups); and \lean{minsize\_unbalanced\_lift},
\lean{minsize\_balanced\_all\_lift}, \lean{minsize\_balanced\_lift}, which are the three theorems above on
\lean{ZMod (p \^{} n * q * r)} with $r\ne2$ prime, \lean{levelsEW w = Es} and \lean{levelsFW w = Fs} written out level by level, and one clause added, that no $\Phi_d$ with $d\mid p^nqr$, $d\nmid p^nq$ divides the mask.
Three more, \lean{Pnq/Sign.lean}, \lean{Pnq/Lift2.lean} and \lean{Pnq/LiftAll.lean} ($1{,}153$ lines), hold the prime $2$ and every $s$:
\lean{Pnq.Lift.exists\_signs} (Lemma~\ref{lem:signs}; a signing is given as a splitting $w=u+v$ of the multiset into the
elements signed $+$ and those signed $-$, and the signed mask is \lean{maskW u - maskW v});
\lean{exists\_lift\_two} (Lemma~\ref{lem:lift2}, for $N=2^k(p^aq^b)$) and \lean{sizes\_eq\_two} (the sizes, as in \lean{sizes\_eq});
\lean{sizes\_eq\_odd} (for every $N$ and every odd $s$, the sizes of the nonempty multisets with exactly a pattern $P$ of divisors of $N$ are the same on
\lean{ZMod (N * s)} and on \lean{ZMod N});
\lean{sizes\_eq\_mul} (the first sentence of Corollary~\ref{cor:lift}: the same on \lean{ZMod (p \^{} n * q * s)} and
\lean{ZMod (p \^{} n * q)}, for primes $p\ne q$ and every $s\ge1$); and \lean{minsize\_unbalanced\_mul},
\lean{minsize\_balanced\_all\_mul}, \lean{minsize\_balanced\_mul}, the three theorems above on
\lean{ZMod (p \^{} n * q * s)} for every $s\ge1$, with the same clause added.
Not formalised: Proposition~\ref{prop:setlift}, Corollary~\ref{cor:sets} and the remarks after them, on sets; the statement on signs for $n=p\,r\,q^b$; and the remarks of Section~\ref{sec:lift} that stand outside its lemmas and Corollary~\ref{cor:lift}: that for odd $N$ Lemmas~\ref{lem:signs} and~\ref{lem:lift2} are one statement, the remark on $r=p$, and the remark after Corollary~\ref{cor:lift} (the three theorems assert their lower bound for the multisets that satisfy the added clause). That lifting along $2$ holds wherever the statement on signs holds for the odd part is in \lean{Pnq/Lift2.lean} as a lemma of the proving instance, \lean{TwoAux.exists\_lift\_two\_aux}, which was not among the statements fixed or audited.

All six files were made in the same way as the eleven. For the first three the author fixed the one definition, one instance and the seven statements first; a fresh instance wrote the
proofs; the author checked that the fixed declarations stand character for character, that Lean prints the
statements as it did before the proofs existed, that there is no \lean{sorry}, added axiom, instance, notation
or attribute beyond the fixed ones, and that the seven theorems rest on the same three axioms. Another fresh
instance audited the statements and found no difference from what is meant; with programs of its own it found
Lean's value to be the least size in 317 of 372 cases in groups up to $\Z_{72}$ and did not finish the other 55
(none came out different), and for $r=2$, which those statements exclude, it found a lift with no new divisor for
every multiset it tried. The three files are in the shared development (\lean{lean-tiling/Pnq/}), where they were put after Noon had read the seven statements; the author then built a fresh copy of the development with them from nothing ($18{,}459$ jobs, no error; Mathlib itself was not rebuilt), and Noon did the same in a copy of its own, with the same count of jobs, the three check files printing $95$ lines that each name only the three axioms. Noon also tried \lean{exists\_lift} by brute force, with nothing from the proof, on $1{,}481$ small multisets for $r=3$ and $r=5$: each had a lift.
For the other three the author fixed the eight statements (they need no new definition) and posted a hash of them on the board that the instances share before any proof was
begun; three fresh instances then wrote the proofs at the same time, one file each; the author made the same checks, with the same result. Another fresh instance audited the eight statements and found no difference from what is meant. With programs of its own it found, by exhaustion, a good signing for each of $1{,}581{,}519{,}831$ multisets on $40$ moduli $p^aq^b$ up to $375$, and a lift along $2$ for each of $696{,}029{,}875$ multisets on $112$ groups up to $\Z_{225}$; equal sets of sizes, as far as its censuses reach, for every pattern in $270$ pairs of groups ($44$, $139$ and $196$ for the three statements, which overlap); and no least size different from Lean's value in $2{,}344$ cases with $s>1$: in $1{,}481$ of them it proved the bound ($1{,}219$ by complete enumeration, $262$ with a solver), and in the other $863$ it found a multiset of Lean's size and none smaller as far as it searched. Noon read the eight statements, compared them with the finished files by a script of its own, and tried \lean{exists\_signs} and \lean{exists\_lift\_two} by exhaustion over small multisets on small groups; the three files were then put into the shared development, and the author built a fresh copy of it with them from nothing ($18{,}465$ jobs, no error; Mathlib itself was not rebuilt). Noon has not yet built them in a copy of its own.

\subsection*{What the formal proof does not show} That the theorem is new; that the prose proofs of this
paper are right (they were read by the author and, in their working-note form, refereed by fresh instances,
which is a weaker thing); anything about sets; and nothing has been read by a person.


\section{Computations}\label{sec:computations}

The formulas of (a) and (b) were written down before, and the rule of (c) was found from, exhaustive
computation, which confirms all three. For each pattern the least size was sought by two separately written models (CP-SAT on class counts,
and HiGHS on polynomial remainders modulo $\Phi_d$), every witness being checked a third way by complex
values: all $2^{2n+1}$ patterns for $n=2$ and $(p,q)\in\{(2,3),(3,2),(2,5),(5,2),(3,5),(5,3)\}$, for
$n=3$ and $(2,3)$, $(3,2)$, $(2,5)$, and for $n=4$ and $(2,3)$; and for $n=3$ and $(3,5)$, $(5,2)$,
$(5,3)$ the $125$, $42$ and $10$ patterns of $128$ that a solver finished: $1{,}265$ patterns decided
in all, $1{,}181$ of them by both solvers and $84$ by the floating-point solver alone. Theorem~\ref{thm:main} gives the computed value in every
one. The rule of (c) was fitted on $n\le3$ and sealed before the $n=4$ results were opened; an earlier guess
of the author's, without clauses (2) and (4), was wrong on $39$ of $554$ patterns. The floating-point solver
is a verdict and not a proof; the formal proof makes it unnecessary, and the table is now a check on the
formal \emph{statement}: Lean's own evaluation of the formula agrees with all $1{,}265$ computed values
(Section~\ref{sec:lean}).

\begin{thebibliography}{9}
\bibitem{CM} E.~Coven and A.~Meyerowitz, \emph{Tiling the integers with translates of one finite set},
 J.~Algebra 212 (1999), 161--174.
\bibitem{KLaMS} G.~Kiss, I.~{\L}aba, C.~Marshall and G.~Somlai, \emph{Lower bounds for mask polynomials with many
cyclotomic divisors}, Adv. Math. 494 (2026), 110932; arXiv:2507.11672. (Read in part by the author; read
 through by a fresh instance of the model.)
\bibitem{KMSV} G.~Kiss, R.~D.~Malikiosis, G.~Somlai and M.~Vizer, \emph{On the discrete Fuglede and Pompeiu
 problems}, Anal. PDE 13 (2020), 765--788; arXiv:1807.02844. (Proposition~3.8 and Theorem~4.9 were read by a
 fresh instance of the model; the author has not read the paper.)
\bibitem{LL} T.~Y.~Lam and K.~H.~Leung, \emph{On vanishing sums of roots of unity}, J.~Algebra 224 (2000),
 91--109; arXiv:math/9511209. (Read in part by a fresh instance of the model; not by the author.)
\bibitem{LaM} I.~{\L}aba and C.~Marshall, \emph{Vanishing sums of roots of unity and the Favard length of
 self-similar product sets}, Discrete Analysis 2022:19; arXiv:2202.07555. (Theorem~1.2, Corollary~6.4 and
 their proofs were read by fresh instances of the model; the author has not read the paper.)
\bibitem{Mal} R.~D.~Malikiosis, \emph{On the structure of spectral and tiling subsets of cyclic groups}, Forum
 Math. Sigma 10 (2022), e23; arXiv:2005.05800. (The proof of Theorem~6.1 was read by a fresh instance of the
 model and by Noon; the author checked the lines derived from it.)
\bibitem{MK} R.~D.~Malikiosis and M.~N.~Kolountzakis, \emph{Fuglede's conjecture on cyclic groups of order
 $p^nq$}, Discrete Analysis 2017:12; arXiv:1612.01328. (The author read pp.~10--11.)
\end{thebibliography}

\end{document}
